\chapter{Plan de Trabajo, Estructura del Proyecto y Presupuesto}
\label{ch:plan_trabajo_presupuesto}

La viabilidad operativa y el éxito de un proyecto de ingeniería avanzada no solo dependen de su solidez científica, sino también de una planificación estratégica rigurosa y un control financiero estricto. En este capítulo se detalla la gobernanza ingenieril que rige el desarrollo del proyecto. En primer lugar, se expone el enfoque metodológico adaptado para la investigación cuántica y se desglosa el pipeline mediante una Estructura de Descomposición del Trabajo (EDT). Posteriormente, se presenta el cronograma temporal del proyecto a través de un diagrama de Gantt, determinando las dependencias de las tareas y el camino crítico de la investigación. Asimismo, se define la estructura organizativa de los recursos humanos involucrados y la asignación de sus responsabilidades técnicas. Finalmente, se detalla el estudio económico, consolidando un presupuesto exhaustivo que contempla los costes de personal, los costes operativos de infraestructura en la nube y cómputo cuántico variacional por shot, las amortizaciones de activos y las licencias de software necesarias para garantizar la reproducibilidad integral del estudio.

\section{Fases del Proyecto y Metodología (EDT)}
\label{sec:fases_proyecto_edt}

\subsection{Enfoque Metodológico}
Dada la naturaleza híbrida de este proyecto, que combina la investigación científica teórica con el desarrollo ágil de software e ingeniería de circuitos cuánticos, se ha adoptado una metodología ágil basada en el marco de trabajo \textit{Scrum}, adaptada de forma específica a entornos de I+D. El ciclo de vida del proyecto se ha estructurado en \textit{Sprints} de dos semanas de duración fija. Al cierre de cada iteración, se han ejecutado pruebas de validación matemática y computacional sobre el repositorio, asegurando que los incrementos de código cuántico mantuvieran la coherencia algebraico-financiera antes de progresar hacia configuraciones experimentales densas.

\subsection{Estructura de Descomposición del Trabajo (EDT)}
La segmentación modular de los objetivos operativos del TFM se ha articulado mediante una Estructura de Descomposición del Trabajo (EDT/WBS), la cual fragmenta el alcance global en cinco Paquetes de Trabajo (\textit{Work Packages}, WP) independientes pero interconectados de forma secuencial:
\begin{itemize}
    \item \textbf{WP1: Gestión y Revisión Bibliográfica:} Comprende la delimitación formal del problema de optimización financiera de carteras, la recopilación del estado del arte en algoritmos variacionales cuánticos y dispositivos NISQ, y el estudio de la clase de complejidad combinatoria asociada a las restricciones operativas reales.
    \item \textbf{WP2: Modelado Matemático y Formulación:} Centrado en el desarrollo algebraico del pipeline de transformación. Integra el diseño de la formulación cuadrática continua de Markowitz, su mapeo sin restricciones a la matriz QUBO con regularizaciones y penalizaciones dinámicas, y la deducción analítica de los coeficientes del Hamiltoniano de Ising.
    \item \textbf{WP3: Desarrollo de Software e Implementación:} Abarca la codificación de la infraestructura nuclear del repositorio. Incluye el módulo de ingesta y cálculo de estadísticas financieras, el diseño del \textit{ansatz} variacional cuántico (preparación del estado de Dicke y la transpilación geométrica del operador mezclador XY) y la integración del bucle de optimización clásica acelerado por JAX.
    \item \textbf{WP4: Experimentación y Benchmarking:} Se enfoca en la ejecución del pipeline experimental. Contempla la simulación de las 40 instancias matemáticas parametrizadas, la recolección homogeneizada de registros en formato tabular y la ejecución comparativa frente a solucionadores clásicos exactos (\textit{Gurobi Optimizer}) y estocásticos (\textit{Simulated Annealing}).
    \item \textbf{WP5: Redacción de la Memoria y Diseminación:} Consiste en el procesado científico de los resultados, la generación automática de figuras analíticas y tablas vectoriales en LaTeX, la redacción formal de la memoria bajo las directrices académicas institucionales y la preparación del material de soporte para la defensa de la investigación.
\end{itemize}

La disposición jerárquica de la EDT y sus subdivisiones en tareas elementales se ilustra en el diagrama estructural de la Figura \ref{fig:wbs_tree}.

\begin{figure}[htbp]
    \centering
    \includegraphics[width=0.85\textwidth]{output/figures/wbs_tree.png}
    \caption{Estructura de Descomposición del Trabajo (EDT/WBS) del proyecto, detallando la jerarquía de los cinco paquetes de trabajo primarios y sus actividades de soporte técnico.}
    \label{fig:wbs_tree}
\end{figure}

\section{Planificación Temporal y Calendario (Gantt)}
\label{sec:planificacion_temporal_gantt}

\subsection{Cronograma de Ejecución e Hitos}
El desarrollo del proyecto se ha proyectado sobre un calendario de ejecución temporal riguroso, adaptado a los plazos académicos de la convocatoria universitaria. La estimación del esfuerzo y la secuenciación de las tareas se ha calibrado para absorber las desviaciones computacionales lógicas derivadas del entrenamiento variacional y la transpilación de circuitos cuánticos complejos. Con el fin de evaluar de forma objetiva la evolución del pipeline del proyecto, se han definido los siguientes Hitos (\textit{Milestones}) clave en el cronograma:
\begin{itemize}
    \item \textbf{M1 (Cierre del WP1):} Estado del arte consolidado y bases del problema de carteras de inversión validadas.
    \item \textbf{M2 (Cierre del WP2):} Formulación analítica QUBO y mapeo de Ising completados con equivalencia matemática certificada.
    \item \textbf{M3 (Cierre del WP3):} Repositorio de software funcional con cobertura de tests automatizados $\ge 80\%$.
    \item \textbf{M4 (Cierre del WP4):} Campaña experimental completada y registros de simulación homogeneizados en almacenamiento estructurado.
    \item \textbf{M5 (Cierre del WP5):} Memoria definitiva del TFM redactada y aprobada para su depósito oficial.
\end{itemize}

\subsection{Determinación del Camino Crítico}
La interdependencia secuencial de los paquetes de trabajo se gestiona mediante la modelización del camino crítico del proyecto. Debido a la naturaleza estrictamente acumulativa de las fases científicas, el camino crítico está integrado de forma rígida por la cadena de tareas correspondiente a los paquetes \textbf{WP2 $\rightarrow$ WP3 $\rightarrow$ WP4}. Cualquier retraso temporal en la formalización del álgebra del QUBO posterga la implementación de los circuitos en Qrisp y, de forma paralela, compromete la ventana de tiempo reservada para las simulaciones variacionales en Google Colab Pro. El cronograma de barras horizontales unificado, que denota las dependencias y resalta el camino crítico operativo, se presenta en la Figura \ref{fig:gantt_chart}.

\begin{figure}[htbp]
    \centering
    \includegraphics[width=0.95\textwidth]{output/figures/gantt_chart.png}
    \caption{Diagrama de Gantt del proyecto, explicitando la duración temporal de los sprints, las dependencias de precedencia analítica y las tareas que integran el camino crítico operativo.}
    \label{fig:gantt_chart}
\end{figure}

\section{Estructura Organizativa y Gestión de Recursos Humanos}
\label{sec:estructura_organizativa_recursos}

\subsection{Definición de Roles y Responsabilidades Técnicas}
Para simular el despliegue del proyecto bajo estándares profesionales de la industria de las finanzas cuantitativas cuánticas (\textit{Quantum Fintech}), la estructura organizativa y de recursos humanos se divide en dos perfiles especializados diferenciados que cubren las competencias clave de la investigación:
\begin{enumerate}
    \item \textbf{Quantum Computing Engineer (Ingeniero Cuántico):} Responsable del diseño físico y algorítmico de los circuitos lógicos. Su esfuerzo se centra en la definición teórica del Hamiltoniano de Ising, la programación de la rutina de preparación del estado de Dicke en Qrisp, la transpilación geométrica sobre la topología lineal LNN y la parametrización de los optimizadores cuánticos variacionales híbridos. Carga de trabajo asignada: $65$ horas laborables.
    \item \textbf{Quantitative Financial Analyst (Analista Cuantitativo):} Responsable del modelado y validación del dominio financiero. Sus tareas comprenden la extracción y preprocesado estadístico de las series temporales de Yahoo Finance, la calibración analítica de la matriz de covarianza, la sintonización de los coeficientes de penalización adaptativos y la evaluación posterior del Ratio de Sharpe y la volatilidad real de las soluciones. Carga de trabajo asignada: $35$ horas laborables.
\end{enumerate}

\subsection{Matriz de Asignación de Responsabilidades (RACI)}
La gobernanza operacional y el reparto de tareas entre los perfiles profesionales se articula mediante la definición de una Matriz RACI (\textit{Responsible, Accountable, Consulted, Informed}). Esta matriz asegura que cada actividad del proyecto cuente con una asignación clara de responsabilidades, tal como se detalla en la Tabla \ref{tab:matriz_raci}.

\begin{table}[htbp]
    \centering
    \caption{Matriz RACI de asignación de responsabilidades por paquete de trabajo.}
    \label{tab:matriz_raci}
    \begin{tabular}{lcc}
        \toprule
        \textbf{Paquete de Trabajo / Actividad EDT} & \textbf{Quantum Computing Engineer} & \textbf{Quantitative Financial Analyst} \\
        \midrule
        \textbf{WP1:} Revisión bibliográfica y estado del arte & A / R & R \\
        \textbf{WP2:} Formulación matemática e Ising & A / R & R \\
        \textbf{WP3:} Desarrollo del repositorio de software & A / R & C \\
        \textbf{WP4:} Experimentación, simulación y benchmarking & R & A / R \\
        \textbf{WP5:} Redacción científica y material didáctico & A / R & A / R \\
        \bottomrule
    \end{tabular}
    \begin{quote}
        \footnotesize \textit{Nomenclatura: A = Accountable (Aprobador de la tarea), R = Responsible (Ejecutor), C = Consulted (Perfil consultado), I = Informed (Perfil informado).}
    \end{quote}
\end{table}

\section{Presupuesto Detallado y Plan de Financiación}
\label{sec:presupuesto_plan_financiacion}

\subsection{Costes Directos de Personal}
El cálculo del coste de capital humano se deduce aplicando tarifas salariales de mercado competitivas, adecuadas a la alta especialización requerida en los ámbitos de la física cuántica y las finanzas algorítmicas. Para el \textit{Quantum Computing Engineer}, se establece un coste por hora de $45$~\euro{}, lo que equivale a un coste bruto total de $2.925$~\euro{} por las $65$ horas de dedicación asignadas. Para el \textit{Quantitative Financial Analyst}, se aplica una tarifa horaria de $25$~\euro{}, consolidando un coste directo de personal de $875$~\euro{} por sus $35$ horas de esfuerzo computado. El montante consolidado por el concepto de recursos humanos asciende de forma estricta a $3.800$~\euro{}.

\subsection{Costes de Infraestructura Cloud y Cómputo Cuántico}
Debido a la inviabilidad física de ejecutar simulaciones cuánticas densas de vectores de estado por encima de los $N \ge 14$ cúbits de forma exclusiva en estaciones locales, el entorno experimental se apoya en recursos distribuidos en la nube. Los gastos de infraestructura informática se calculan bajo una modalidad de pago por uso (\textit{Pay-as-you-go}) y se desglosan en las siguientes partidas analíticas:
\begin{itemize}
    \item \textbf{Cómputo Clásico Acelerado (Google Colab Pro / AWS EC2):} Despliegue de instancias con aceleración de hardware por GPU (NVIDIA A100/V100) dedicadas al motor de diferenciación automática JAX y el simulador JaspQAOA, garantizando la compilación Just-In-Time de los circuitos variacionales. Gasto estimado: $400$~\euro{}.
    \item \textbf{Acceso a Unidades de Procesamiento Cuántico (AWS Braket):} Coste derivado del envío de ráfagas operativas a procesadores NISQ reales de compuertas lógicas (como los dispositivos basados en superconductores de Rigetti o IQM) para testear canales de ruido controlados. El gasto de facturación se calcula considerando una tasa fija de apertura de tarea por circuito ($0.30$~\euro{}) sumada a una tasa por disparo o muestra individual ($0.0002$~\euro{} por \textit{shot}), asumiendo la escala estándar de $1024$ disparos fijada por la especificación. Gasto estimado: $2.000$~\euro{}.
\end{itemize}

La suma consolidada del aprovisionamiento en la nube e infraestructura cloud asciende a un total de $2.400$~\euro{}.

\subsection{Licencias de Software y Amortización de Activos}
Para asegurar el rigor analítico e industrial del benchmarking clásico, se contempla la adquisición de una licencia comercial del entorno de optimización matemática lineal y cuadrática entera mixta \textit{Gurobi Optimizer} prorrateada de forma exacta según la ventana temporal de uso efectivo del proyecto, lo que arroja un coste de licencia de $1.500$~\euro{}. Por otra parte, la amortización técnica de la estación de trabajo local empleada para la fase de desarrollo básico de código y ejecución de los tests unitarios automáticos se calcula bajo una vida útil de 36 meses, repercutiendo un valor prorrateado de $400$~\euro{} al coste de desarrollo de la investigación.

\subsection{Presupuesto Financiero Consolidado}
Complementando los conceptos computados con una partida de contingencias prudencial del $10\%$ para amortiguar posibles imprevistos operativos en el consumo de recursos de cómputo, se consolida el cuadro macroeconómico del proyecto en la Tabla \ref{tab:presupuesto_total}, reflejando un coste integral neto de $8.910$~\euro{}.

\begin{table}[htbp]
    \centering
    \caption{Presupuesto general consolidado del proyecto de investigación.}
    \label{tab:presupuesto_total}
    \begin{tabular}{lp{6.2cm}cc}
        \toprule
        \textbf{Categoría Económica} & \textbf{Descripción Técnico-Económica de la Partida} & \textbf{Coste Unitario} & \textbf{Coste Total} \\
        \midrule
        \textbf{Personal (RRHH)} & Quantum Computing Engineer ($65$ h) & $45$ \euro{}/h & $\hphantom{0}2.925$ \euro{} \\
        & Quantitative Financial Analyst ($35$ h) & $25$ \euro{}/h & $\hphantom{00}875$ \euro{} \\
        \addlinespace
        \textbf{Infraestructura Cloud} & Instancias clásicas GPU AWS EC2 / Colab & Pago por uso & $\hphantom{00}400$ \euro{} \\
        & Tiempo neto de QPU por \textit{shot} AWS Braket & Tarifa por disparo & $\hphantom{0}2.000$ \euro{} \\
        \addlinespace
        \textbf{Software Licenses} & Suscripción Gurobi Optimization API & Tarifa plana & $\hphantom{0}1.500$ \euro{} \\
        \addlinespace
        \textbf{Amortizaciones} & Estación de trabajo y hardware local & Ciclo 36 meses & $\hphantom{00}400$ \euro{} \\
        \addlinespace
        \textbf{Contingencias} & Margen de desviación operativa ($10\%$) & — & $\hphantom{00}810$ \euro{} \\
        \midrule
        \textbf{TOTAL (IVA excl.)} & & & \textbf{\hphantom{0}8.910 \euro{}} \\
        \bottomrule
    \end{tabular}
\end{table}